
\documentclass{article}
\usepackage{colortbl}
\usepackage{makecell}
\usepackage{multirow}
\usepackage{supertabular}

\begin{document}

\newcounter{utterance}

\centering \large Interaction Transcript for game `cladder', experiment `full\_v1.5\_default', episode 8830 with refineonly{-}Qwen3.5{-}9B.
\vspace{24pt}

{ \footnotesize  \setcounter{utterance}{1}
\setlength{\tabcolsep}{0pt}
\begin{supertabular}{c@{$\;$}|p{.15\linewidth}@{}p{.15\linewidth}p{.15\linewidth}p{.15\linewidth}p{.15\linewidth}p{.15\linewidth}}
    \# & \multicolumn{2}{c}{Player} && \multicolumn{2}{c}{Game Master} \\
    \hline

    \theutterance \stepcounter{utterance}  
    & & & \multicolumn{4}{p{0.6\linewidth}}{
        \cellcolor[rgb]{0.9,0.9,0.9}{
            \makecell[{{p{\linewidth}}}]{
                \texttt{\tiny{[P1$\langle$GM]}}
                \texttt{You are an expert in causal inference. The following question is not a typical commonsense query, but rather a meticulously designed question created by a professor specializing in causal inference, intended to assess the students' mastery of the course content.} \\
\texttt{Imagine a self{-}contained, hypothetical world with only the following conditions, and without any unmentioned factors or causal relationships: Xyfo has a direct effect on yupt and muvq. Zupj has a direct effect on yupt. Yupt has a direct effect on muvq. Xyfo is unobserved. Method 1: We look at how yupt correlates with muvq case by case according to xyfo. Method 2: We look directly at how yupt correlates with muvq in general. To understand how yupt affects muvq, is it more correct to use the Method 1 than Method 2?} \\
\texttt{Start your answer with 'Yes' or 'No', followed by additional reasoning or evidence to support your explanation.} \\
            }
        }
    }
     \\ \\

    \theutterance \stepcounter{utterance}  
    & \multicolumn{4}{p{0.6\linewidth}}{
        \cellcolor[rgb]{0.561,0.737,0.561}{
            \makecell[{{p{\linewidth}}}]{
                \texttt{\tiny{[P1$\rangle$GM]}}
                \texttt{No} \\
\\ 
\texttt{**Reasoning:**} \\
\\ 
\texttt{In this hypothetical world, the causal structure can be represented as follows:} \\
\texttt{1.  \$X\_\{yfo\} \textbackslash{}rightarrow Y\_\{upt\}\$} \\
\texttt{2.  \$X\_\{yfo\} \textbackslash{}rightarrow M\_\{uvq\}\$} \\
\texttt{3.  \$Z\_\{upj\} \textbackslash{}rightarrow Y\_\{upt\}\$} \\
\texttt{4.  \$Y\_\{upt\} \textbackslash{}rightarrow M\_\{uvq\}\$} \\
\\ 
\texttt{The goal is to estimate the causal effect of \$Y\_\{upt\}\$ on \$M\_\{uvq\}\$.} \\
\\ 
\texttt{**Analysis of the Confounding:**} \\
\texttt{There is a backdoor path between \$Y\_\{upt\}\$ and \$M\_\{uvq\}\$: \$Y\_\{upt\} \textbackslash{}leftarrow X\_\{yfo\} \textbackslash{}rightarrow M\_\{uvq\}\$.} \\
\texttt{Because \$X\_\{yfo\}\$ is an unobserved common cause (a confounder) of both the treatment (\$Y\_\{upt\}\$) and the outcome (\$M\_\{uvq\}\$), the marginal correlation between \$Y\_\{upt\}\$ and \$M\_\{uvq\}\$ will be biased. This bias arises because \$X\_\{yfo\}\$ influences both variables simultaneously, creating a spurious association that does not reflect the true causal mechanism of \$Y\_\{upt\} \textbackslash{}rightarrow M\_\{uvq\}\$.} \\
\\ 
\texttt{**Evaluation of Method 2:**} \\
\texttt{Method 2 looks at the general correlation between \$Y\_\{upt\}\$ and \$M\_\{uvq\}\$ without conditioning on anything. Since \$X\_\{yfo\}\$ is unobserved, we cannot control for it. Therefore, Method 2 yields a biased estimate of the causal effect due to the confounding by \$X\_\{yfo\}\$.} \\
\\ 
\texttt{**Evaluation of Method 1:**} \\
\texttt{Method 1 looks at the correlation between \$Y\_\{upt\}\$ and \$M\_\{uvq\}\$ "case by case according to \$X\_\{yfo\}\$." This phrasing implies stratifying or conditioning on the levels of \$X\_\{yfo\}\$ (i.e., comparing the relationship within groups where \$X\_\{yfo\}\$ is held constant).} \\
\texttt{By conditioning on \$X\_\{yfo\}\$, we block the backdoor path \$Y\_\{upt\} \textbackslash{}leftarrow X\_\{yfo\} \textbackslash{}rightarrow M\_\{uvq\}\$. Within each stratum of \$X\_\{yfo\}\$, the variation in \$Y\_\{upt\}\$ is no longer correlated with the variation in \$M\_\{uvq\}\$ via the confounder. Consequently, the association observed within these strata isolates the direct causal effect of \$Y\_\{upt\}\$ on \$M\_\{uvq\}\$ (assuming no other unmeasured confounders exist, which the prompt guarantees by stating "only the following conditions").} \\
\\ 
\texttt{**Conclusion:**} \\
\texttt{Since Method 2 fails to address the unobserved confounding by \$X\_\{yfo\}\$, while Method 1 successfully controls for this confounder through stratification, Method 1 provides the correct causal estimate. Therefore,} \\
            }
        }
    }
    & & \\ \\

    \theutterance \stepcounter{utterance}  
    & & & \multicolumn{2}{p{0.3\linewidth}}{
        \cellcolor[rgb]{0.9,0.9,0.9}{
            \makecell[{{p{\linewidth}}}]{
                \texttt{\tiny{[GM$|$GM]}}
                \texttt{No} \\
            }
        }
    }
    & & \\ \\

    \theutterance \stepcounter{utterance}  
    & & & \multicolumn{2}{p{0.3\linewidth}}{
        \cellcolor[rgb]{0.9,0.9,0.9}{
            \makecell[{{p{\linewidth}}}]{
                \texttt{\tiny{[GM$|$GM]}}
                \texttt{yes} \\
            }
        }
    }
    & & \\ \\

    \theutterance \stepcounter{utterance}  
    & & & \multicolumn{2}{p{0.3\linewidth}}{
        \cellcolor[rgb]{0.9,0.9,0.9}{
            \makecell[{{p{\linewidth}}}]{
                \texttt{\tiny{[GM$|$GM]}}
                \texttt{game\_result = LOSE} \\
            }
        }
    }
    & & \\ \\

\end{supertabular}
}

\end{document}
